\subsubsection{Giải chi tiết Bài tập c (Phương pháp đơn hình)}

\textbf{1. Đưa bài toán về dạng chuẩn tắc}

Các ràng buộc chính của bài toán đều có dạng "$\le$", do đó ta chỉ cần cộng thêm vào vế trái của mỗi ràng buộc một biến phụ không âm $x_4, x_5, x_6 \ge 0$ để biến chúng thành đẳng thức.
Mô hình bài toán dạng chuẩn được viết lại như sau:

\begin{alignat*}{4}
    f(x) = \quad & -2x_1 & {}+{} & 5x_2 & {}+{} & 3x_3 & {}+{} 0x_4 + 0x_5 + 0x_6 \rightarrow \max \\
    \text{s.t.} \quad & x_1 & {}+{} & 2x_2 & {}+{} & 3x_3 & {}+{} x_4 = 8000 \\
    & x_1 & {}+{} & 3x_2 & {}+{} & x_3 & {}+{} x_5 = 10000 \\
    & x_1 & {}+{} & 2x_2 & {}+{} & x_3 & {}+{} x_6 = 4000 \\
    & \rlap{$x_j \ge 0, \quad j = \overline{1,6}$}
\end{alignat*}

Ta có hệ vector điều kiện của các biến phụ $x_4, x_5, x_6$ tạo thành ma trận đơn vị $I_3$, ta chọn tập biến $(x_4, x_5, x_6)$ làm cơ sở xuất phát.

\vspace{0.3cm}
\textbf{2. Quá trình giải bằng Bảng đơn hình}

\textbf{Bảng 1 (Bước lặp 0):}
\begin{table}[H]
    \centering
    \renewcommand{\arraystretch}{1.3}
    \begin{tabular}{|c|c|c|c|c|c|c|c|c|c|}
        \hline
        \multirow{2}{*}{\textbf{Cơ sở}} & \multirow{2}{*}{$\mathbf{C_b}$} & \multirow{2}{*}{\textbf{PA}} & $x_1$ & $x_2$ & $x_3$ & $x_4$ & $x_5$ & $x_6$ & \multirow{2}{*}{\textbf{Tỷ số $\lambda$}} \\
        \cline{4-9}
        & & & -2 & 5 & 3 & 0 & 0 & 0 & \\
        \hline
        $x_4$ & 0 & 8000 & 1 & 2 & 3 & 1 & 0 & 0 & $8000/2 = 4000$ \\
        \hline
        $x_5$ & 0 & 10000 & 1 & 3 & 1 & 0 & 1 & 0 & $10000/3 \approx 3333.3$ \\
        \hline
        $x_6$ & 0 & 4000 & 1 & \textbf{2} & 1 & 0 & 0 & 1 & $4000/2 = \mathbf{2000} \rightarrow$ \\
        \hline
        \multicolumn{2}{|c|}{$\Delta_j$} & \textbf{Z = 0} & 2 & \textbf{-5} $\uparrow$ & -3 & 0 & 0 & 0 & \\
        \hline
    \end{tabular}
\end{table}
\textit{Giải thích xoay bảng 1:}
\begin{itemize}[itemsep=0pt, parsep=0pt, topsep=0pt]
    \item Xét dòng $\Delta_j$, ta có $\Delta_2 = -5$ là giá trị âm nhỏ nhất $\Rightarrow$ Chọn $x_2$ đưa \textbf{vào cơ sở}.
    \item Lập tỷ số $\lambda_i$, ta có $\min\{4000, 3333.3, 2000\} = 2000$ tại dòng $x_6$ $\Rightarrow$ Chọn $x_6$ đưa \textbf{ra khỏi cơ sở}. Phần tử trụ là $\mathbf{2}$.
\end{itemize}

\vspace{0.3cm}
\textbf{Bảng 2 (Bước lặp 1):} Xoay bảng theo phần tử trụ $2$ bằng phép biến đổi Gauss-Jordan.
\begin{table}[H]
    \centering
    \renewcommand{\arraystretch}{1.3}
    \begin{tabular}{|c|c|c|c|c|c|c|c|c|c|}
        \hline
        \multirow{2}{*}{\textbf{Cơ sở}} & \multirow{2}{*}{$\mathbf{C_b}$} & \multirow{2}{*}{\textbf{PA}} & $x_1$ & $x_2$ & $x_3$ & $x_4$ & $x_5$ & $x_6$ & \multirow{2}{*}{\textbf{Tỷ số $\lambda$}} \\
        \cline{4-9}
        & & & -2 & 5 & 3 & 0 & 0 & 0 & \\
        \hline
        $x_4$ & 0 & 4000 & 0 & 0 & \textbf{2} & 1 & 0 & -1 & $4000 / 2 = \mathbf{2000} \rightarrow$ \\
        \hline
        $x_5$ & 0 & 4000 & -1/2 & 0 & -1/2 & 0 & 1 & -3/2 & - \\
        \hline
        $x_2$ & 5 & 2000 & 1/2 & 1 & 1/2 & 0 & 0 & 1/2 & $2000 / (1/2) = 4000$ \\
        \hline
        \multicolumn{2}{|c|}{$\Delta_j$} & \textbf{Z = 10000} & 9/2 & 0 & \textbf{-1/2} $\uparrow$ & 0 & 0 & 5/2 & \\
        \hline
    \end{tabular}
\end{table}
\textit{Giải thích xoay bảng 2:}
\begin{itemize}[itemsep=0pt, parsep=0pt, topsep=0pt]
    \item Bảng vẫn còn $\Delta_3 = -1/2 < 0$, do đó phương án hiện hành chưa tối ưu $\Rightarrow$ Chọn $x_3$ đưa \textbf{vào cơ sở}.
    \item Lập tỷ số $\lambda$ ở cột $x_3$ (lưu ý bỏ qua dòng 2 vì có hệ số $a_{23} = -1/2 \le 0$): ta có $\min\{4000/2, 2000/(1/2)\} = 2000$ $\Rightarrow$ Chọn $x_4$ đưa \textbf{ra khỏi cơ sở}. Phần tử trụ là $\mathbf{2}$.
\end{itemize}

\vspace{0.3cm}
\textbf{Bảng 3 (Bước lặp 2):} Cập nhật bảng theo phần tử trụ $2$.
\begin{table}[H]
    \centering
    \renewcommand{\arraystretch}{1.3}
    \begin{tabular}{|c|c|c|c|c|c|c|c|c|}
        \hline
        \multirow{2}{*}{\textbf{Cơ sở}} & \multirow{2}{*}{$\mathbf{C_b}$} & \multirow{2}{*}{\textbf{PA}} & $x_1$ & $x_2$ & $x_3$ & $x_4$ & $x_5$ & $x_6$ \\
        \cline{4-9}
        & & & -2 & 5 & 3 & 0 & 0 & 0 \\
        \hline
        $x_3$ & 3 & 2000 & 0 & 0 & 1 & 1/2 & 0 & -1/2 \\
        \hline
        $x_5$ & 0 & 5000 & -1/2 & 0 & 0 & 1/4 & 1 & -7/4 \\
        \hline
        $x_2$ & 5 & 1000 & 1/2 & 1 & 0 & -1/4 & 0 & 3/4 \\
        \hline
        \multicolumn{2}{|c|}{$\Delta_j$} & \textbf{Z = 11000} & 9/2 & 0 & 0 & 1/4 & 0 & 9/4 \\
        \hline
    \end{tabular}
\end{table}

\textbf{3. Kết luận}

Tại Bảng 3, tất cả các hệ số ước lượng đều không âm ($\Delta_j \ge 0, \forall j$), do đó phương án hiện hành là phương án tối ưu thỏa mãn tiêu chuẩn dừng.
\begin{itemize}
    \item Phương án tối ưu của bài toán là: $x^* = (0, 1000, 2000)$.
    \item Giá trị lớn nhất của hàm mục tiêu là: $f(x)_{\max} = 11000$.
\end{itemize}